===== BEGIN FILE: sidispt.tex =====
\documentclass[prl,aps,twocolumn,superscriptaddress]{revtex4}
\usepackage{psfrag,slashed,cancel,array,graphicx,todonotes,hyperref}


\usepackage[utf8]{inputenc}
\usepackage{mathtools}
\usepackage{mathrsfs}
\usepackage{multirow}
\usepackage{braket}
\usepackage{slashed}
\usepackage{booktabs}
\usepackage[normalem]{ulem}
\usepackage{slashed}
\hypersetup{pdftitle={},pdfcreator={},linkcolor=[rgb]{0.15,0.35,0.75},colorlinks=true,citecolor=[rgb]{0.675,0,0.2},urlcolor=[rgb]{0.15,0.35,0.65}}



\newcommand{\img}{\mathrm{i}}
\newcommand{\sdt}{\!\cdot\!}
\newcommand{\tr}{\textrm{tr}}
\newcommand{\fr}{\frac}
\newcommand{\ra}{\rightarrow}
\newcommand{\lra}{\leftrightarrow}




\def\beq{\begin{equation}}
\def\eeq{\end{equation}}
\def\bsp#1\esp{\begin{split}#1\end{split}}
\def\d{{\rm d}}
\newcommand{\refcite}[1]{ref.~\cite{#1}}
\newcommand{\refscite}[1]{refs.~\cite{#1}}
\newcommand{\Eq}[1]{Eq.~\eqref{eq:#1}}
\newcommand{\Eqs}[2]{Eqs.~\eqref{eq:#1} and \eqref{eq:#2}}
\newcommand{\eq}[1]{eq.~\eqref{eq:#1}}
\newcommand{\eqs}[2]{eqs.~\eqref{eq:#1} and \eqref{eq:#2}}
\renewcommand{\sec}[1]{Section~\ref{sec:#1}}
\newcommand{\secs}[2]{sections~\ref{sec:#1} and \ref{sec:#2}}
\newcommand{\subsec}[1]{section~\ref{subsec:#1}}
\newcommand{\app}[1]{appendix~\ref{app:#1}}
\newcommand{\fig}[1]{figure~\ref{fig:#1}}
\newcommand{\figs}[2]{figures~\ref{fig:#1} and \ref{fig:#2}}
\newcommand{\cI}{\mathcal{I}}
\newcommand{\cK}{\mathcal{K}}
\newcommand{\cJ}{\mathcal{J}}
\newcommand{\cL}{\mathcal{L}}
\newcommand{\cM}{\mathcal{M}}
\newcommand{\cN}{\mathcal{N}}
\newcommand{\cO}{\mathcal{O}}
\newcommand{\GammaC}{\Gamma_{\rm cusp}}
\newcommand{\nn}{\nonumber}
\newcommand{\MSbar}{\overline{\mathrm{MS}}}
\newcommand{\lqcd}{\Lambda_\mathrm{QCD}}
\newcommand{\born}{ \hat \sigma_0}
\newcommand{\df}{\mathrm{d}}
\newcommand{\bn}{{\bar n}}
\newcommand{\bq}{{\bar q}}
\newcommand{\Sl}[1]{\slashed{#1}}
\newcommand{\SCETII}{SCET$_{\text{II}}$ }
\newcommand{\as}{\alpha_s}
\newcommand{\bt}{{\vec b}_T}
\newcommand{\qt}{{\vec q}_T}
\newcommand{\blue}[1]{{\color{blue}{#1}}}
\newcommand{\gammathr}{\gamma_{\rm th}}
\newcommand{\SYM}{$\mathcal{N}=4$~sYM\,}
\newcommand{\R}{R}
\newcommand{\nslash}{\Sl{n}}
\newcommand{\bnslash}{\Sl{\bn}}
\newcommand{\red}[1]{{\color{red}{#1}}}
\newcommand{\FierzSplit}{{\color{red}\Big|}}
\newcommand{\orange}[1]{{\color{orange}{#1}}}


\AtBeginDocument{
\heavyrulewidth=.08em
\lightrulewidth=.05em
\cmidrulewidth=.03em
\belowrulesep=.65ex
\belowbottomsep=0pt
\aboverulesep=.4ex
\abovetopsep=0pt
\cmidrulesep=\doublerulesep
\cmidrulekern=.5em
\defaultaddspace=.5em
}

\setcounter{topnumber}{1}
\setcounter{tocdepth}{2}

\def\cA{\mathcal{A}}
\def\cM{\mathcal{M}}
\def\cP{\mathcal{P}}
\def\cO{\mathcal{O}}
\def\cC{\mathcal{C}}
\def\cN{\mathcal{N}}
\def\cV{\mathcal{V}}
\def\cT{\mathcal{T}}
\def\cR{\mathcal{R}}

\def\eps{\epsilon}


\def\be{\begin{equation}}
\def\ee{\end{equation}}




\newcommand{\shao}[1]{\textcolor[rgb]{0,0,1}{[Shao: #1]}}
\newcommand{\fang}[1]{\textcolor[rgb]{0,0,1}{[Fang: #1]}}
\begin{document}

\preprint{XXX}


\author{Liang Dong}
\email{liang.dong@sjtu.edu.cn}
\affiliation{State Key Laboratory of Dark Matter Physics, Shanghai Key Laboratory for Particle Physics and Cosmology, Key Laboratory for Particle Astrophysics and Cosmology (MOE),
School of Physics and Astronomy, Shanghai Jiao Tong University, Shanghai 200240, China}
\author{Shen Fang}
\email{sfang23@m.fudan.edu.cn}
\affiliation{Department of Physics, Center for Field Theory and Particle Physics, and Key Laboratory of Nuclear Physics and Ion-beam Application (MOE), Fudan University, Shanghai, 200433, China}
\author{Jun Gao}
\email{jung49@sjtu.edu.cn}
\affiliation{State Key Laboratory of Dark Matter Physics, Shanghai Key Laboratory for Particle Physics and Cosmology, Key Laboratory for Particle Astrophysics and Cosmology (MOE),
School of Physics and Astronomy, Shanghai Jiao Tong University, Shanghai 200240, China}
\author{Hai Tao Li}
\email{haitao.li@sdu.edu.cn}
\affiliation{School of Physics, Shandong University, Jinan, Shandong 250100, China}
\author{Ding Yu Shao}
\email{dyshao@fudan.edu.cn}
\affiliation{Department of Physics, Center for Field Theory and Particle Physics, and Key Laboratory of Nuclear Physics and Ion-beam Application (MOE), Fudan University, Shanghai, 200433, China}
\affiliation{Shanghai Research Center for Theoretical Nuclear Physics, NSFC and Fudan University, Shanghai 200438, China}
\affiliation{Center for High Energy Physics, Peking University, Beijing 100871, China}
\author{Yu Jiao Zhu}
\email{yzhu@mpp.mpg.de}
\affiliation{Max-Planck-Institut fur Physik, Werner-Heisenberg-Institut, Boltzmannstr. 8, 85748 Garching, Germany}

\title{NNLO QCD corrections to hadron production in DIS at finite transverse momentum}

\begin{abstract}

We present the \textit{first} complete calculation of hadron production in deep-inelastic scattering (DIS) at finite transverse momentum to next-to-next-to-leading order (NNLO) in perturbative QCD.

To overcome the long-standing challenge of infrared divergences in semi-inclusive processes with identified final state hadrons at finite transverse momentum, we implement the recently developed $q_T$-subtraction framework based on the recoil-free jet definition.

By utilizing the winner-take-all recombination scheme, we achieve a consistent factorization for hadron-jet associated production, enabling the inclusion of $\mathcal{O}(\alpha_s^3)$ corrections.

Our results demonstrate a significantly improved stabilization of the perturbative expansion and a reduction in scale uncertainties compared to previous next-to-leading order predictions.

We find that the NNLO corrections are essential for a robust description of high precision multiplicity data from the ZEUS collaborations.

This work provides a high precision theoretical foundation for the upcoming Electron-Ion Collider era and establishes a new benchmark for the exploration of the nucleon's three-dimensional structure.

\end{abstract}



\maketitle

\paragraph*{Introduction.---}

The semi-inclusive deep inelastic scattering (SIDIS) process, $l(l) + N(P) \to l'(l') + h(P_h) + X$, is fundamental to mapping the three-dimensional nucleon structure.

While inclusive DIS has historically provided a precise mapping of collinear parton distribution functions (PDFs), the detection of a final state hadron with finite transverse momentum $P_{hT}$ enables the exploration of the transverse momentum dependent (TMD) dynamics of partons and the hadronization process \cite{Collins:1981uk, Collins:1981uw}.

In particular, the high-$P_{hT}$ region is of great theoretical interest.

It provides a unique window into the transition between the non-perturbative TMD framework and the perturbative collinear factorization regime, where the $P_{hT}$ is generated by hard radiation.



Over the past decades, experimental programs have provided a detailed mapping of $P_{hT}$ distributions across diverse kinematic ranges.

Early measurements by the H1 \cite{H1:1996muf, H1:1999iry, H1:2004xgw} and ZEUS \cite{ZEUS:1995acw, ZEUS:1999ies} collaborations at HERA established the initial benchmarks for inclusive hadron production at high $P_{hT}$.

These were followed by higher-precision studies from HERMES \cite{HERMES:2012uyd}, COMPASS \cite{COMPASS:2013bfs, COMPASS:2017mvk}, and Jefferson Lab (JLab) \cite{JeffersonLabHallA:2016ctn}, which revealed intricate flavor dependencies and spin-momentum correlations within the nucleon.



Despite these experimental advances, theoretical predictions for the transverse momentum spectrum have largely been restricted to next-to-leading order (NLO) accuracy \cite{Kniehl:2004hf, Daleo:2004pn, Wang:2019bvb}, which exhibit significant scale uncertainties and struggle to provide a simultaneous global description of experimental data \cite{Gonzalez-Hernandez:2018ipj}.

This status contrasts sharply with rapid progress for inclusive quantities in DIS: N$^3$LO corrections have been computed for single jet production in DIS \cite{Currie:2018fgr}, as well as NNLO corrections for dijet production in DIS \cite{Currie:2017tpe}.

Furthermore, NNLO coefficient functions were recently derived for inclusive SIDIS cross sections~\cite{Abele:2021nyo, Goyal:2023zdi, Bonino:2024qbh, Bonino:2024wgg, Goyal:2024tmo, Goyal:2024emo, Bonino:2025tnf, Haug:2025ava, Bonino:2025qta, Zhou:2025lqv, Bonino:2025bqa, Goyal:2025qyu, Gao:2026tnd}.



The fully differential description of identified hadrons at finite $P_{hT}$ remains a critical frontier, and achieving NNLO precision is essential to stabilize the perturbative expansion and provide robust theory inputs for the upcoming Electron-Ion Collider (EIC) \cite{Accardi:2012qut, AbdulKhalek:2021gbh}.

However, extending calculations to this order for final states involving identified particles faces significant hurdles due to the complex cancellation of infrared (IR) divergences \cite{Kinoshita:1962ur, Lee:1964is}.

For the description of inclusive quantities, various strategies have emerged to resolve these singularities, broadly categorized into local subtraction \cite{Gehrmann-DeRidder:2005btv, Czakon:2010td, Currie:2013vh}, projection-to-Born methods \cite{Cacciari:2015jma}, and slicing techniques \cite{Catani:2007vq, Gao:2012ja, Gao:2014eea, Gaunt:2015pea, Boughezal:2015dva, Berger:2016inr}.

Yet, applying these techniques to semi-inclusive processes has been limited by the difficulty of consistently treating soft and collinear radiation in the presence of identified hadrons and jets, with a few successful applications~\cite{Fu:2024fgj, Czakon:2025yti, Generet:2025bqx, Gao:2026tnd, Bonino:2026dvr}.


In this Letter, we present the \textit{first} complete NNLO QCD calculation of identified hadron production in DIS at finite $P_{hT}$.

The core of our theoretical approach is the implementation of a new $q_T$-subtraction framework recently proposed in Ref.~\cite{Fu:2024fgj, Fu:2026nkd} (see also Refs.~\cite{Buonocore:2022mle, Buonocore:2023rdw, Buonocore:2025ucn}), which generalizes the classical transverse momentum subtraction method~\cite{Catani:2007vq} to final states involving both jets and hadrons.

By utilizing a recoil-free jet axis, specifically the winner-take-all (WTA) recombination scheme~\cite{Salam:WTAUnpublished, Bertolini:2013iqa}, this method allows for the consistent treatment of radiation from the identified hadron and other jets.

By incorporating $\mathcal{O}(\alpha_s^3)$ corrections through this framework, we demonstrate good convergence of theoretical predictions and reduced theoretical uncertainties.

This result marks a new era of precision for SIDIS phenomenology and provides a rigorous test of the QCD factorization theorems in the high-$P_{hT}$ domain.


\begin{figure}[t!]
    \centering
    \includegraphics[width=0.4\textwidth]{DIS_2jet.pdf}
    \caption{Schematic definition of the kinematic variables used for the $q_T$-subtraction method in the Breit frame. The incoming virtual photon $\gamma^*$ scatters off the proton, producing an identified hadron and a leading recoiling jet ${\cal J}_k$. The blue lines along the proton direction denote beam collinear radiation. The green curved lines indicate the soft radiation.  The slicing variable is defined by the azimuthal decorrelation $\delta\phi$ (or equivalently the out-of-plane momentum $p_{\rm out}$ of the hadron with respect to the beam-jet plane). The limit $\delta\phi \to 0$ corresponds to the unresolved back-to-back configuration governed by the factorization formula in Eq.~(\ref{eq:fact_phi_jet}).}
\label{fig:kinematics}
\end{figure}

\
\paragraph*{$q_T$-subtraction formalism.---}

In this work, we employ the azimuthal decorrelation $\delta\phi$, or equivalently the hadron momentum component $p_{\rm out}$ perpendicular to the scattering plane, as the slicing variable.

The scattering kinematics in the Breit frame are illustrated in Fig.~\ref{fig:kinematics} with the scattering plane defined by the leading recoiling jet and beam directions.

This partitions the cross section into contributions from an unresolved region, dominated by soft and collinear radiation, and a resolved region characterized by large-angle emissions.

Explicitly, we have
\begin{align}
 {\d\sigma\over \d\mathcal O}= \!\!\int_0^{\delta \phi^{\rm cut}}\d \delta\phi \frac{\d\sigma}{\d\delta \phi\,\d\mathcal O} +\!\!\int_{\delta \phi^{\rm cut}}^{\delta \phi^{\rm max}} \d \delta\phi \frac{\d\sigma}{ \d\delta \phi\,\d\mathcal O},
\end{align}

where $\delta \phi^{\rm cut}$ is a small resolution/slicing parameter introduced in the $q_T$-subtraction method.

The first (unresolved) part is approximated by the leading-power (LP) contribution in the unresolved limit ($\delta\phi^{\rm cut}\to0$), subject to power corrections which vanish in this limit.

The cancellation of the dependence on $\delta\phi^{\rm cut}$ in the full results serves as a rigorous check of IR safety and numerical stability of our framework.




In the resolved region ($\delta\phi > \delta \phi^{\rm cut}$) where the hadron and jets are well separated, the differential cross sections can be calculated by requiring production of an inclusive hadron in association with a dijet in DIS up to NLO.

We use the one-loop matrix elements for $ei \to ejkl$ and the tree-level matrix elements for $ei \to ejklm$ calculated in Refs.~\cite{Bern:1997sc, Campbell:2002tg, Campbell:2003hd, Campbell:2010ff}, where $ijklm$ represent all possible flavor combinations of QCD partons.

The jets are reconstructed using the $k_T$ algorithm~\cite{Catani:1993hr, Ellis:1993tq} with WTA recombination scheme.

Crucially, the jet algorithm dependence cancels between the resolved and unresolved regions, ensuring the final result is algorithm independent.

The calculations are carried out in the {\tt FMNLO} framework~\cite{Liu:2023fsq} for identified hadron production, with IR singularities handled using a method of dipole subtraction~\cite{Catani:1996vz} in combination with phase space slicing.


\begin{figure*}[t!]
    \centering
    \includegraphics[width=0.95\textwidth]{afig1.pdf}
    \caption{
    Numerical stability of the total cross section with respect to the slicing parameter $\delta\phi^{\rm cut}$.

    The panels display the $\mathcal{O}(\alpha_s^2)$ (red) and $\mathcal{O}(\alpha_s^3)$ (blue) contributions for six representative partonic channels.

    Error bars indicate Monte Carlo statistical uncertainties.

    The solid cyan lines represent the fit to the asymptotic limit $\delta\phi^{\rm cut} \to 0$.

    The dashed purple lines denote the reference values of $\mathcal{O}(\alpha_s^2)$.}
    \label{fig:cutoff_dep}
\end{figure*}


In the unresolved region ($\delta\phi < \delta \phi^{\rm cut}$), the LP cross section is governed by a factorization formula in impact parameter ($b$) space, which captures the azimuthal decorrelation between the identified hadron and the recoiling jet in the back-to-back limit.

A central feature of our framework is the WTA recombination scheme.

Because the WTA axis is recoil-free, it remains insensitive to soft radiation, providing a stable reference for the hard scattering direction.

This property eliminates non-global logarithms \cite{Dasgupta:2001sh, Neill:2016vbi}, substantially simplifying the structure of the factorization formula \cite{Gutierrez-Reyes:2018qez, Chien:2020hzh}.

Building upon the soft-collinear effective theory (SCET) framework~\cite{Bauer:2000yr, Bauer:2001ct, Bauer:2001yt, Bauer:2002nz, Beneke:2002ph} and following recent developments in Refs.~\cite{Chien:2020hzh, Chien:2022wiq, Fang:2023thw, Fang:2024auf, Fu:2024fgj}, we obtain the factorized cross section
  \begin{align} \label{eq:fact_phi_jet}
      &\frac{\d\sigma_{\text{LP}}}{\d x \, \d y \, \d z \, \d^2 \vec P_{hT}  \, \d p_{\rm out}} = \int \frac{\d b}{2\pi}e^{i p_{\rm out} b/\zeta} \\
  &\hspace{1cm}\times   \sum_{ijk} \int \d \xi\,   H_{ei\to ejk}(Q,\xi,\zeta)\nn\\
  &\hspace{1cm}\times  {\cal B}_{i/p}(\xi,b)\,
      {\cal D}_{h/j}(\zeta,b)\,{\cal J}_{k}(b)  \, S_{ij k}(b)\,, \nn
\end{align}
where $p_{\rm out}$ is the hadron's transverse momentum relative to the scattering plane.

The standard kinematic variables are defined as
\begin{align}
Q^2=-q^2\,,~ x=\frac{Q^2}{2P\cdot q}\,, ~ y=\frac{P\cdot q}{P\cdot l}\,,~z=\frac{P\cdot P_h}{P\cdot q}\,,
\end{align}
with virtual photon momentum $q=l-l'$.

The factorized ingredients $H$, $\cal B$, $\cal D$, $S$, and $\mathcal{J}$ represent the hard, TMD beam, fragmentation, soft, and WTA jet functions, respectively, where dependence on the renormalization and rapidity scales $(\mu,\nu)$ is implicit. $\xi$ and $\zeta=z+xP_{hT}^2/[z (\xi-x) Q^2]$ are the momentum fractions for the beam and fragmentation functions.

Finally, the sum $\sum_{ijk}$ runs over all contributing partonic channels, specifically photon-gluon fusion ($\gamma^* g\to q \bar q$) and QCD Compton scattering ($\gamma^* q\to q g$).



Achieving NNLO accuracy requires two-loop corrections for each factorized ingredient.

The hard function $H_{ei\to ejk}$ encodes the tree-level and virtual corrections to the partonic process $ei \to ejk$.

We extract them from~\cite{Garland:2002ak, Gehrmann:2009vu} and find full agreement with those used in Refs.~\cite{Gehrmann:2022vuk, Gehrmann:2023zpz, NNLOJET:2025rno}.

Crucially, linearly polarized contributions from both incoming and outgoing gluons are essential for a complete description of the cross section beyond LO \cite{Catani:2010pd, Chien:2020hzh,  Gao:2023ivm}.



We employ standard TMD beam ($\mathcal{B}_{i/p}$) and fragmentation ($\mathcal{D}_{h/j}$) functions~\cite{Collins:2011zzd, Boussarie:2023izj}, applying their well-established NNLO matching onto collinear distribution functions $f_{a/p}$ and $d_{h/b}$ \cite{Catani:2011kr, Catani:2012qa, Gehrmann:2012ze, Gehrmann:2014yya, Luebbert:2016itl, Echevarria:2016scs, Luo:2019hmp, Luo:2019bmw}
\begin{align}
  {\cal B}_{i/p}\left(\xi,b\right) &=\sum_a \int_{\xi}^1 \frac{\d {\hat\xi}}{\hat\xi} \, {\cal I}_{i / a}(\hat\xi, b ) \,  f_{a/p}\left(\frac{\xi}{\hat\xi}, \mu\right)\,,   \nonumber \\
  {\cal D}_{h/j}\left(\zeta,b\right) &=
\frac{1}{\zeta^2}\sum_b \int_{\zeta}^1 \frac{\d {\hat\zeta }}{\hat\zeta} \, {\cal C}_{b / j}(\hat\zeta, b ) \, d_{h/b}\left(\frac{\zeta}{\hat\zeta}, \mu\right) \,,
\end{align}
with ${\cal I}_{i / a}=\delta_{ia}\delta(1-\hat\xi)+\mathcal{O}(\alpha_s)$ and ${\cal C}_{b / j}=\delta_{bj}\delta(1-\hat\zeta)+\mathcal{O}(\alpha_s)$.

The soft function $S_{ijk}$ follows the universal NNLO TMD soft function definition \cite{Echevarria:2015byo, Lubbert:2016rku}, retaining a simple analytic form because the hard scattering process is in a planar configuration  \cite{Gao:2019ojf, Chien:2020hzh, Chien:2022wiq, Gao:2023ivm, Fu:2024fgj}.



Furthermore, we adopt the $k_T$ algorithm and the WTA recombination scheme for the TMD jet functions $\mathcal{J}_k$, to be consistent with the numerical evaluation in the resolved region.

The NNLO quark jet functions are readily available from recent calculations \cite{Gutierrez-Reyes:2019vbx, Fang:2024auf, Bell:2021dpb, Brune:2022cgr, Buonocore:2025ysd}.

While the NNLO WTA gluon jet functions, comprising unpolarized and linearly polarized components, are currently unknown, renormalization invariance uniquely determines the two-loop logarithmic terms.

We approximate the remaining two constants of the gluon jet functions with rescaled energy-energy correlation jet functions \cite{Luo:2019hmp,Luo:2019bmw}, and find that their numerical impact is negligible.



\paragraph*{Numerical results.---}

To validate our calculation framework, we explicitly verify that the physical cross section remains independent of the arbitrary cutoff $\delta\phi^{\rm cut}$ in the limit $\delta\phi^{\rm cut}\to 0$.

Fig.~\ref{fig:cutoff_dep} displays the numerical stability of the $\mathcal{O}(\alpha_s^3)$ (NNLO) and $\mathcal{O}(\alpha_s^2)$ (NLO) corrections to the differential cross section of charged pion production at the EIC.

The kinematics are fixed at $\sqrt{s}=141$ GeV, $x=0.2$, $y=0.3$, with integration over $0.2<z<0.8$ and $P_{hT} > z Q/10$.

The results are shown for six representative partonic channels, including $uu$, $ug$, $gu$, $gg$, $u\bar u$, and $u d$, labeled according to the flavors of the incoming and fragmenting partons.

Among these, the $uu$ channel is always dominant due to both favored PDFs and FFs.

We use the CT18 NNLO PDFs~\cite{Hou:2019efy} and the NPC23 NNLO FFs~\cite{Gao:2025hlm} with a nominal choice of $\mu=Q$ for both the renormalization and factorization scales.

The NLO results (red points) exhibit excellent flatness over the range $10^{-3} < \delta\phi^{\rm cut} < 10^{-1}$, confirming a precise cancellation of IR divergences, and converge to the reference values calculated with {\tt FMNLO}~\cite{Liu:2023fsq}.

At NNLO (blue points), the results similarly converge to a constant plateau, despite the larger Monte Carlo integration uncertainties from the resolved cross section.

The solid lines represent a fit to the genuine NNLO corrections together with residual power corrections~\cite{Grazzini:2017mhc, Ebert:2019zkb}.

From the fit, we estimate a systematic uncertainty (discrepancy \textit{w.r.t} the genuine result) of less than 5\%(2\%) when a finite value of $\delta\phi^{\rm cut}=0.04(0.01)$ is used in the NNLO calculations for all partonic channels.

The stability observed across both quark- and gluon-initiated channels establishes the robustness of our method for high precision phenomenology.



\begin{figure}[t!]
    \includegraphics[width=0.9\linewidth]{afig2.pdf}
    \caption{
    Differential cross section as functions of the hadron momentum fraction $z$ for charged pion production at the EIC at various orders in QCD together with scale variations.

    Kinematic cuts, including a lower limit of 2~GeV on the hadron transverse momentum, are applied.

    }
    \label{fig:z_dist}
\end{figure}



Fig.~\ref{fig:z_dist} presents our NNLO predictions for the charged pion differential cross section as a function of the momentum fraction $z$, tailored to the kinematics of the future EIC, including a transverse momentum cut of $P_{hT} > 2$~GeV.

The upper panel displays the perturbative progression from LO (blue dotted) to NLO (red dot-dashed) and NNLO (black solid).

We observe a substantial positive correction at NNLO across the entire $z$ range, highlighting the importance of the $\mathcal{O}(\alpha_s^3)$ contributions.

The lower panel, which normalizes results to the central LO prediction, reveals the dramatic stabilization of the perturbative expansion.

The scale uncertainty bands (hatched regions), calculated by varying $\mu$ by a factor of two from the nominal value $Q$,  undergo a significant reduction when moving from NLO to NNLO, confirming that the missing higher-order terms are now under control.

Furthermore, the $K$-factor exhibits a distinct kinematic dependence, rising noticeably at large $z$.

This behavior indicates that NNLO corrections are particularly crucial in the threshold region, where soft-gluon dynamics begin to dominate~\cite{Abele:2022wuy, Goyal:2025bzf}.




\begin{figure}[t!]
    \includegraphics[width=0.9\linewidth]{afig3.pdf}
    \caption{
    Differential cross section as functions of the hadron transverse momentum $P_{hT}$ for charged pion production at the EIC at various orders in QCD together with scale variations.

    }
    \label{fig:pt_dist}
\end{figure}



Fig.~\ref{fig:pt_dist} displays the differential cross section as a function of the hadron transverse momentum $P_{hT}$, which is of particular significance as it directly probes the transition from the non-perturbative TMD regime to the perturbative collinear domain.

The upper panel reveals a clear hierarchy in the perturbative series, with the NNLO prediction (black solid line) providing a refined description of the spectrum over the entire kinematic range.

The ratio plot in the lower panel demonstrates the excellent convergence of the expansion.

While the jump from LO to NLO is substantial—reflecting the opening of new partonic channels—the correction from NLO to NNLO is moderate and stable.

Crucially, the scale uncertainty (hatched bands) is largely reduced at NNLO compared to the lower orders.

This improved theoretical precision is essential for reliably identifying the matching window between fixed-order calculations and TMD resummation, thereby facilitating a robust extraction of the nucleon's three-dimensional structure at the EIC.



\begin{figure}[t!]
    \includegraphics[width=0.9\linewidth]{afig4.pdf}
    \caption{
    Multiplicity distribution as functions of the hadron transverse momentum $P_{hT}$ for unidentified charged hadron production at the HERA at various orders in QCD together with scale variations, compared to the ZEUS data~\cite{ZEUS:1995acw}.

    }
    \label{fig:hera}
\end{figure}



We compare our predictions with high-precision measurements from the ZEUS collaboration \cite{ZEUS:1995acw} in Fig.~\ref{fig:hera}.

The plot displays the unidentified charged hadron multiplicity as a function of transverse momentum $P_{hT}$ in the kinematic region with $10< Q^2 <160$ GeV$^2$, $75< W <175$ GeV, and $x_F>0.05$, where the Feynman $x$ variable is defined as $x_F\equiv 2P_{hL}^*/W$.

The longitudinal momentum $P_{hL}^*$ is the projection of hadron momentum onto the virtual photon direction in the center of mass frame of $\gamma^*p$, and $W$ is the invariant mass of the hadronic final state.

The LO result (blue dotted) grossly underestimates the data, accounting for only $\sim 50\%$ of the observed yield. This large discrepancy highlights the limitation of LO kinematics in generating high-$P_{T}$ recoil.

While NLO predictions (green) significantly reduce this gap, they continue to systematically undershoot the experimental central values and show scale variations as large as the experimental errors.

The inclusion of $\mathcal{O}(\alpha_s^3)$ corrections enhances the normalization for large transverse momentum, and exhibits improved convergence together with reduced scale variations.

This agreement demonstrates that NNLO precision is indispensable for the quantitative interpretation of SIDIS data at high energies.


\paragraph*{Summary and Outlook.---}

In this Letter, we have presented the first complete NNLO QCD calculation for single inclusive hadron production in DIS at finite transverse momentum.

By implementing the $q_T$-subtraction method with a recoil-free WTA axis, we have successfully overcome the technical bottleneck of consistently treating IR divergences in the presence of both fragmenting hadrons and jets.

Our numerical results demonstrate that NNLO corrections are substantial, playing a critical role in stabilizing the perturbative series and reducing scale uncertainties,
and provide a robust quantitative description of the hadron spectrum in the high energy regime.



Our results initiate a new precision frontier for SIDIS phenomenology.

By reliably fixing the collinear cross section, we can precisely delineate the transition region between collinear and TMD factorization, facilitating a deeper understanding of the nucleon's three-dimensional structure.

Future extensions of this work to polarized scattering will further unlock the potential of the EIC to probe the spin-momentum correlations of partons within hadrons.

Finally, our results also serve as a fundamental step toward fully differential N$^3$LO predictions for SIDIS~\cite{Gao:2026tnd}.





\paragraph*{Acknowledgments.---}

We thank Guido Bell, Kevin Brune, Xuan Chen, Thomas Gehrmann, Ian Moult, and Hua Xing Zhu for helpful discussions.

The work of LD and JG is supported by the National Natural Science Foundation of China (NSFC) under Grant No. 12275173 and the Shanghai Municipal Education Commission under Grant No. 2024AIZD007.

SF and DYS are supported by the National Science Foundations of China under Grant No.~12275052, No.~12147101, No.~12547102, and the Innovation Program for Quantum Science and Technology under grant No. 2024ZD0300101.

HTL is supported by the National Science Foundation of China under Grant Nos. 12275156 and 12321005.

The computations in this paper were run on the Siyuan-1 cluster supported by the Center for High Performance Computing at Shanghai Jiao Tong University.


\bibliographystyle{apsrev4-2}
\bibliography{sidispt.bib}


\end{document}
===== END FILE: sidispt.tex =====
